% ARGUES: the measurement pipeline is sound and the analysis was fixed in advance.
% DRAWS FROM: PRE_REGISTRATION_OSF.md §4, §5, §6, §8.
% STATUS: drafted 2026-08-12. Counts through verification are final. Judge-side
%         counts move until all 31,980 citations are judged.

\section{Methods}
\label{sec:methods}

\subsection{Question set}

The question set has 900 items, 100 in each cell of a three-domain by
three-cell design. It was frozen and hashed before any model was queried, and
the hash is published in the pre-registration, so the set cannot have been
adjusted after seeing how models performed.

The three cells differ in what an honest answer looks like. \emph{Known-real}
items ask for citations that exist and are well indexed, recruited half from a
hyper-canonical pool and half from a moderately-known pool, so canonicality can
be tested as a covariate instead of assumed irrelevant. \emph{Obscure-but-real}
items ask for work that exists but sits in the under-cited tail, recruited at a
citation count in $[2, 50)$. That band was calibrated to the empirical knowledge
inflection; an arbitrary low threshold would have put the cell below the range
where the distinction we need is meaningful.
\emph{Unanswerable} items have no correct citation at all.

The unanswerable cell carries a construction burden the other two do not,
because an item is only unanswerable if the thing it asks about does not exist.
Each candidate passed a two-step negative-existence proof: a Levenshtein-3 fuzzy
title match against Crossref, OpenAlex, PubMed and CourtListener, rejecting on
any partial match, followed by a SPECTER v1 semantic similarity check rejecting
any candidate within cosine 0.85 of a top-5 title-search hit. Calibration on 30
known near-duplicate pairs rejected 30 of 30, with the lowest known-duplicate
cosine at 0.912. The same threshold rejects 3 of 6 related-but-different paper
pairs, so the check over-rejects rather than under-rejects. For a
negative-existence proof that is the safe direction, and it means the accepted
unanswerable items are more conservatively novel than a strict same-paper
detector would produce.

Known-real and obscure-but-real items were restricted to pre-2023 publications,
the intersection of the model training cutoffs. This filter was applied
inconsistently at first: the legal known-real recruiter filtered on citation
count without a year gate, leaving 45 opinions dated 2023 or later in a cell
whose whole purpose is to hold capability constant. Post-cutoff cases in a
capability-ceiling control introduce a recency confound. All 45 were replaced
with same-tier pre-2023 opinions drawn by seeded draw from the same candidate
pool, under the amendment procedure the pre-registration specifies for this
case. The amended file is operative for every reported analysis, the original is
retained unmodified with its own hash, and the two can be diffed.

Domain scopes are narrower than the domain labels suggest, and we describe them
narrowly throughout. \emph{Legal} means US federal case law: Supreme Court,
Circuit Courts of Appeals, and reported district-court opinions. State courts
are excluded because CourtListener's coverage of unreported state decisions is
incomplete, which would put accidental collisions into the unanswerable cell.
\emph{Academic} means literature indexed under a broad OpenAlex citation-graph
concept, which in practice spans bioinformatics, materials science, psychology,
statistics and core machine learning rather than machine learning alone.
\emph{Medical} means a pool that is roughly 60\% Cochrane reviews, which
over-represents oncology, cardiology, infectious disease and the
systematic-review genre.

\begin{table}[htbp]
\centering
\small
\caption{Composition of the frozen 900-question set. The design is balanced by
construction: every domain contributes exactly 100 questions to every cell. The
strata inside known-real are the pre-registered canonicality tiers, which exist
so that recall accuracy can be checked against canonicality rather than assumed
to be uniform. Both real cells are restricted to pre-2023 publications.}
\label{tab:dataset}
\begin{tabular}{@{}l l r l@{}}
\toprule
Cell & Domain & $n$ & Recruitment stratum \\
\midrule
Known-real       & Academic & 50  & OpenAlex, $>$9{,}999 citations (hyper-canonical) \\
                 &          & 50  & OpenAlex, 1{,}000--5{,}000 citations (moderate) \\
                 & Legal    & 50  & SCOTUS, $\geq$100 citing references \\
                 &          & 50  & Federal Circuit, $\geq$50 citing references \\
                 & Medical  & 60  & Cochrane Reviews, $>$49 citations \\
                 &          & 40  & Non-Cochrane, $>$499 citations \\
\midrule
Obscure-but-real & Academic & 100 & OpenAlex, 2--49 citations \\
                 & Legal    & 100 & Federal district court, 2--49 citing references \\
                 & Medical  & 100 & OpenAlex, 2--49 citations \\
\midrule
Unanswerable     & Academic & 100 & Constructed; no real referent by design \\
                 & Legal    & 100 & Constructed; no real referent by design \\
                 & Medical  & 100 & Constructed; no real referent by design \\
\midrule
\textbf{Total}   &          & \textbf{900} & \\
\bottomrule
\end{tabular}
\end{table}

Forty-five of the legal known-real questions in this table are the replacements
described above rather than the originally drawn opinions.

\subsection{Models}

The primary tests use open-weights checkpoints, for a structural reason rather
than a practical one. Testing whether post-training raises fabrication requires
a base-model comparator, and closed-API models have none: within that
population, whether a model is post-trained is collinear with which model it is.
The open-weights cascade provides base, supervised-fine-tuned,
preference-optimized and reinforcement-tuned checkpoints of the same base
models, so \texttt{stage\_ordinal} varies within a family and the interaction is
estimable.

Every checkpoint is pinned to a commit SHA and every closed-API call to a dated
snapshot. Where a snapshot was deprecated mid-experiment, both were run and
reported as separate rows rather than substituted.

\begin{table}[htbp]
\centering
\small
\caption{The open-weights checkpoints. Three cascades, each carrying all four
post-training stages of one base model, which is what makes
\texttt{stage\_ordinal} vary within a family and the interaction estimable.
Revisions are the pinned Hugging Face commits actually served.}
\label{tab:models}
\begin{tabular}{@{}l l r l l@{}}
\toprule
Family & Checkpoint & Ord. & Stage & Revision \\
\midrule
\texttt{tulu3\_8b}  & \texttt{meta-llama/Llama-3.1-8B}              & 0 & base     & \texttt{d04e592bb4f6} \\
                    & \texttt{allenai/Llama-3.1-Tulu-3-8B-SFT}      & 1 & SFT      & \texttt{f2a0b46b0cfd} \\
                    & \texttt{allenai/Llama-3.1-Tulu-3-8B-DPO}      & 2 & DPO      & \texttt{a7beb67e33ff} \\
                    & \texttt{allenai/Llama-3.1-Tulu-3-8B}          & 3 & RLVR     & \texttt{666943798adb} \\
\midrule
\texttt{tulu3\_70b} & \texttt{meta-llama/Llama-3.1-70B}             & 0 & base     & \texttt{349b2ddb53ce} \\
                    & \texttt{allenai/Llama-3.1-Tulu-3-70B-SFT}     & 1 & SFT      & \texttt{f8ef34145946} \\
                    & \texttt{allenai/Llama-3.1-Tulu-3-70B-DPO}     & 2 & DPO      & \texttt{2dab5183da92} \\
                    & \texttt{allenai/Llama-3.1-Tulu-3-70B}         & 3 & RLVR     & \texttt{cfc1d855e534} \\
\midrule
\texttt{olmo2\_13b} & \texttt{allenai/OLMo-2-1124-13B}              & 0 & base     & \texttt{3fefddc1bf18} \\
                    & \texttt{allenai/OLMo-2-1124-13B-SFT}          & 1 & SFT      & \texttt{b7ec47ed94f3} \\
                    & \texttt{allenai/OLMo-2-1124-13B-DPO}          & 2 & DPO      & \texttt{97a7f33fc8b4} \\
                    & \texttt{allenai/OLMo-2-1124-13B-Instruct}     & 3 & Instruct & \texttt{3a5c85baefbb} \\
\bottomrule
\end{tabular}
\end{table}

Ordinal 0 is the base checkpoint and \texttt{is\_post\_trained} is 0 there and 1
everywhere else. Ordinal 3 is the final released model of each cascade, which is
RLVR for the two Tülu-3 families and the Instruct release for OLMo-2; the
registered coding treats both as the terminal stage. Two distinct base models
are represented, Llama-3.1 and OLMo-2, and two post-training recipes, Tülu-3's
and OLMo-2's, so the 8B and 70B Tülu-3 cascades share a recipe while differing
in scale.

\subsection{Generation and prompt conditions}

Each model answered every question under four prompt conditions.
\texttt{neutral} asks the question with no further instruction.
\texttt{citations\_expected} asks for at least two academic citations in
support. \texttt{abstention\_permitted} adds that uncertainty is preferred over
a wrong citation and that saying so is better than guessing.
\texttt{length\_matched} asks for a thorough multi-paragraph answer of
comparable length to the \texttt{citations\_expected} prompt while mentioning
neither citations nor sources.

The fourth condition exists to answer an objection that would otherwise sink
E2. A longer answer mechanically affords more citation slots, so a model asked
for citations might fabricate more simply because it wrote more. The
\texttt{citations\_expected} against \texttt{length\_matched} contrast separates
the citation cue from the verbosity demand, and the pre-registration commits in
advance to how each outcome is reported, including the case where the contrast
runs against our hypothesis (Appendix~\ref{app:deviations}).

Open-weights generations were produced at temperature 0, one per model,
question and condition, giving 43{,}200 generations.

\subsection{Citation extraction and verification}

Citations were extracted from generation text by a snapshot-pinned model with
forced tool use, yielding 40{,}840 citations, of which 25{,}530 generations
carried at least one. Base models cite far less often than post-trained ones.
That gap is itself a finding, and it also thins the base-model rows that the E1
contrast depends on, so we report per-model citation counts alongside the
fabrication rates.

Each citation was then checked against Crossref, Semantic Scholar, OpenAlex,
PubMed and CourtListener. Legal citations route differently and this matters:
the academic indexes do not cover US case law, so an academic-index miss on a
legal citation is uninformative. CourtListener is the authoritative existence
check for that domain, and a legal citation it cannot match is recorded as
unverifiable rather than fabricated. Without that rule, real but
academically-unindexed federal opinions, which is what the legal
obscure-but-real cell is made of, would be counted as fabrications by
construction.

Verification returned REAL for 8{,}860 citations (21.7\%), FABRICATED for
17{,}504 (42.9\%) and UNVERIFIABLE for 14{,}476 (35.4\%). The Semantic Scholar
API key was unavailable for these runs. The direction of that loss is
one-sided: a missing source can only turn a real citation unverifiable or
fabricated, never turn a fabricated one real, so it inflates apparent
fabrication and the judge stage corrects it.

\begin{table}[htbp]
\centering
\small
\caption{Pipeline yield. Each of the twelve checkpoints answered all 900
questions under all four conditions, so every cell rests on exactly 3{,}600
generations by construction and no cell is thinned by attrition. ``Citing'' is
the share of generations producing at least one extractable citation.}
\label{tab:t3}
\begin{tabular}{@{}l l r r r r@{}}
\toprule
\multicolumn{6}{@{}l}{\emph{By checkpoint}} \\
Family & Stage & Generations & Citing & Citations & Per gen. \\
\midrule
\texttt{tulu3\_8b}  & base     & 3{,}600 & 21.3\% & 951    & 0.264 \\
                    & SFT      & 3{,}600 & 61.5\% & 3{,}197 & 0.888 \\
                    & DPO      & 3{,}600 & 79.8\% & 4{,}947 & 1.374 \\
                    & RLVR     & 3{,}600 & 80.3\% & 5{,}069 & 1.408 \\
\texttt{tulu3\_70b} & base     & 3{,}600 & 18.7\% & 852    & 0.237 \\
                    & SFT      & 3{,}600 & 66.6\% & 3{,}261 & 0.906 \\
                    & DPO      & 3{,}600 & 84.1\% & 5{,}000 & 1.389 \\
                    & RLVR     & 3{,}600 & 84.2\% & 4{,}997 & 1.388 \\
\texttt{olmo2\_13b} & base     & 3{,}600 & 5.3\%  & 327    & 0.091 \\
                    & SFT      & 3{,}600 & 54.9\% & 2{,}730 & 0.758 \\
                    & DPO      & 3{,}600 & 76.1\% & 4{,}784 & 1.329 \\
                    & Instruct & 3{,}600 & 76.3\% & 4{,}725 & 1.313 \\
\midrule
\multicolumn{6}{@{}l}{\emph{By condition}} \\
\multicolumn{2}{@{}l}{\texttt{neutral}}              & 10{,}800 & 67.7\% & 9{,}414  & 0.872 \\
\multicolumn{2}{@{}l}{\texttt{citations\_expected}}  & 10{,}800 & 66.6\% & 15{,}954 & 1.477 \\
\multicolumn{2}{@{}l}{\texttt{abstention\_permitted}}& 10{,}800 & 31.9\% & 4{,}950  & 0.458 \\
\multicolumn{2}{@{}l}{\texttt{length\_matched}}      & 10{,}800 & 70.2\% & 10{,}522 & 0.974 \\
\midrule
\multicolumn{2}{@{}l}{\textbf{Total}}                & 43{,}200 & 59.1\% & 40{,}840 & 0.945 \\
\bottomrule
\end{tabular}
\end{table}

Two features of Table~\ref{tab:t3} matter for reading the results. First, base
checkpoints barely cite at all: 5.3\% of OLMo-2 base generations produce a
citation, against 76\% for its Instruct release. The base column throughout this
paper is therefore a floor set partly by non-compliance with the instruction, and
not a demonstration that base models are more careful. Second,
\texttt{citations\_expected} does not produce the most \emph{citing generations};
\texttt{length\_matched} does, at 70.2\% against 66.6\%. What
\texttt{citations\_expected} maximizes is citations \emph{per generation}, 1.48
against 0.97. That distinction is what the length-matched control was built to
exploit.

\subsection{The judge}

Verification alone is not accurate enough to carry the outcome variable. It
agrees with blind human labels on 46\% of a hand-checked sample, because it
cannot tell a citation that does not exist from one its sources do not reach,
and because it cannot evaluate whether a real work supports the claim made of
it.

A judge stage resolves those cases. It is an LLM with web search and fetch, and
it fires only on the 31{,}980 citations that verification called FABRICATED or
UNVERIFIABLE. Citations that verification confirmed as real are not re-examined.
The combine rule is stated here exactly because it is easy to paraphrase
wrongly: a verifier REAL is accepted as is; a judge verdict of REAL or
FABRICATED overrides a verifier FABRICATED or UNVERIFIABLE; and a judge verdict
of UNVERIFIABLE falls back to the verifier label rather than rescuing it.

The judge applies a dividing line that the taxonomy needs and that we state as a
rule rather than leaving to case-by-case judgment. A real work cited with an
incorrect year or venue is a metadata error and counts as real. A real title
attached to a different real work's authors is not a metadata error, because the
citation now identifies a work that does not exist, and it counts as fabricated.

The step that accepts verifier REAL calls without review was validated on
earlier data and fails on this corpus. It is corrected by an amendment and the
correction is described in Appendix~\ref{app:deviations}, since it changes which
citations are re-examined rather than how any of them is scored.

\subsection{Statistical specification}

Everything in this subsection was registered before the data existed.

The unit of analysis is the generation, not the citation, so that responses
which cite nothing stay in the model. A generation is coded \texttt{fabricated =
1} if any of its citations is fabricated or misgrounded, and 0 if all of its
citations verify \emph{or} it produced no citation at all. Coding abstention as
0 is deliberate and pre-committed: abstention is honest non-fabrication, and the
reward-hacking account specifically predicts that an honest model abstains on
obscure and unanswerable items, so abstentions must enter the model rather than
disappear from it. Generations that cite only unverifiable sources are excluded
from the primary analysis, which is the one exclusion that conditions on
something downstream of treatment, and Section~\ref{sec:results} reports the
selection check that the pre-registration requires before that exclusion is
interpreted.

Both confirmatory models are mixed-effects logistic regressions fit with
\texttt{lme4::glmer}, with crossed random intercepts for model nested in family
and for question:

\begin{verbatim}
m_e1 <- glmer(fabricated ~ stage_ordinal * cell * domain
                         + (1 | model_family/model_id) + (1 | question_id), ...)

m_e2 <- glmer(fabricated ~ condition * is_post_trained * cell + domain
                         + (1 | model_family/model_id) + (1 | question_id), ...)
\end{verbatim}

The random-effects structure is parsimony-primary by registration, with the
design-justified maximal structure reported as a labeled sensitivity. This
departs from the keep-it-maximal recommendation on purpose, and the reason is
specific to this design: with few families carrying a usable stage contrast, a
family-level random slope is estimated from too few grouping levels to be
stable, and for a binomial model a design-justified slope can absorb the
between-group variance that the fixed interaction would otherwise capture. The
maximal model is therefore not uniformly conservative here, and powering on it
would understate the registered test. Which model supplies the reported p-value
was fixed in advance by an operative-model rule, not chosen after seeing both.

Correction is Holm-Bonferroni within experiment family, not across the paper,
and each family is enumerated in advance with $m = 5$. The E1 family is the
\texttt{stage\_ordinal} main effect, its interactions with the obscure-real and
unanswerable cells, and its three-way interactions with the medical and legal
domains. The E2 family is the \texttt{citations\_expected} and
\texttt{abstention\_permitted} effects, each of their interactions with
\texttt{is\_post\_trained}, and the three-way interaction of
\texttt{citations\_expected}, \texttt{is\_post\_trained} and the obscure-real
cell. The strictest threshold any single test can face is therefore $\alpha/5 =
0.01$, and the power analysis reports a-priori power at that floor as well as at
$\alpha = 0.05$.

Effect sizes are odds ratios with Wald confidence intervals on the target
interaction terms. The registration also asked for parametric-bootstrap
percentile intervals from 10,000 resamples; at a measured cost of about 19,200
worker-hours we did not compute them, and amendment \#14
(Appendix~\ref{app:deviations}) gives the measurement and the reasoning. Fabricated and misgrounded rates are reported
separately as well as summed, because a result driven by citations that do not
exist and one driven by real work cited for claims it does not support carry
different interpretations, even though the confirmatory test sums them.

Of the registered sensitivity analyses, two bear directly on the results below.
One recodes unverifiable citations as fabricated. It excludes no generation at
all, so it dissolves the missingness question instead of assuming an answer to
it; the pre-registration specifies a trigger for elevating it to co-primary, and
Appendix~\ref{app:deviations} records that we elevated it. The other treats the
continuous and factor codings of \texttt{stage\_ordinal} as co-primary under an
intersection-union rule, so the headline claim requires both to hold. An earlier
draft reported whichever coding came out significant, which is two chances at
the same test and inflates the family-wise error rate to roughly 9.75\%.
